\documentclass[11pt]{article}

\usepackage[margin=1in]{geometry}
\usepackage{amsmath}
\usepackage{amssymb}
\usepackage{hyperref}

\title{Stage 3F Physics Note: Conformal Cartoon Translation}
\author{GL-with-AI Project}
\date{\today}

\begin{document}
\maketitle

\section{Scope}

Stage 3F translates the diagonal reduced cartoon geometry from physical metric
variables into the conformal variables used by GRChombo-style CCZ4. It is a
derivation note only. No C++ source terms or variables are implemented.

\section{Metric Dictionary}

For gridded indices \(A,B\in\{x,z\}\),
\begin{equation}
  \gamma_{AB}=\chi^{-1}h_{AB},
  \qquad
  \gamma_{ww}=\chi^{-1}h_{ww}.
\end{equation}
The diagonal physical metric is therefore
\begin{equation}
  dl^2
  =
  \chi^{-1}h_{xx}\,dx^2
  +
  \chi^{-1}h_{zz}\,dz^2
  +
  x^2\chi^{-1}h_{ww}\,d\Omega_2^2.
\end{equation}
Define
\begin{equation}
  q=\gamma_{ww}=\frac{h_{ww}}{\chi}.
\end{equation}
The physical transverse sphere radius is
\begin{equation}
  R_{\rm sphere}=x\sqrt{\gamma_{ww}}
  =
  x\sqrt{\frac{h_{ww}}{\chi}}.
\end{equation}
The angular component is
\begin{equation}
  g_{\theta\theta}=x^2\gamma_{ww}
  =
  \frac{x^2h_{ww}}{\chi},
\end{equation}
not merely \(h_{ww}\) or \(\gamma_{ww}\).

\section{Determinant Condition}

The conformal split is not fully specified until the conformal determinant
condition is imposed. For the diagonal cartoon/Cartesian-like
four-spatial-dimensional metric block,
\begin{equation}
  \det\gamma_{\rm 4D}
  =
  \gamma_{xx}\gamma_{zz}\gamma_{ww}^2,
\end{equation}
because the two hidden directions both contribute \(\gamma_{ww}\). Thus
\begin{equation}
  \chi
  =
  \left(\det\gamma_{\rm 4D}\right)^{-1/4}
  =
  \left(\gamma_{xx}\gamma_{zz}\gamma_{ww}^2\right)^{-1/4},
\end{equation}
and
\begin{equation}
  h_{xx}=\chi\gamma_{xx},
  \qquad
  h_{zz}=\chi\gamma_{zz},
  \qquad
  h_{ww}=\chi\gamma_{ww}.
\end{equation}
The reduced conformal determinant is then
\begin{equation}
  \det h_{\rm 4D}=h_{xx}h_{zz}h_{ww}^2=1.
\end{equation}
This is the cartoon/Cartesian-like determinant condition. It does not include
the spherical-coordinate Jacobian factor \(x^4\sin^2\theta\). The angular
metric component remains
\begin{equation}
  g_{\theta\theta}
  =
  x^2\gamma_{ww}
  =
  \frac{x^2h_{ww}}{\chi}.
\end{equation}

\section{Derivative Identities}

For \(q=h_{ww}/\chi\),
\begin{equation}
  \partial_A q
  =
  \chi^{-1}\partial_A h_{ww}
  -
  \chi^{-2}h_{ww}\partial_A\chi .
\end{equation}
The second derivative is
\begin{align}
  \partial_A\partial_B q
  &=
  \chi^{-1}\partial_A\partial_B h_{ww}
  -
  \chi^{-2}
  \left[
    \partial_A h_{ww}\partial_B\chi
    +
    \partial_B h_{ww}\partial_A\chi
    +
    h_{ww}\partial_A\partial_B\chi
  \right] \nonumber\\
  &\quad
  +
  2\chi^{-3}h_{ww}\partial_A\chi\,\partial_B\chi.
\end{align}
Equivalently,
\begin{equation}
  \partial_A\log q
  =
  \partial_A\log h_{ww}
  -
  \partial_A\log\chi,
\end{equation}
and
\begin{equation}
  \frac{\partial_A\partial_B q}{q}
  =
  \partial_A\partial_B\log q
  +
  \partial_A\log q\,\partial_B\log q.
\end{equation}

\section{Physical Radius}

Let
\begin{equation}
  f=x\sqrt{q}=x\sqrt{\frac{h_{ww}}{\chi}}.
\end{equation}
Then
\begin{align}
  \partial_x\log f
  &=
  \frac{1}{x}
  +
  \frac{1}{2}
  \left(\partial_x\log h_{ww}-\partial_x\log\chi\right),\\
  \partial_z\log f
  &=
  \frac{1}{2}
  \left(\partial_z\log h_{ww}-\partial_z\log\chi\right),
\end{align}
and
\begin{equation}
  \partial_A f=f\,\partial_A\log f.
\end{equation}

\section{Extrinsic Curvature}

With \(\mathrm{GR\_SPACEDIM}=4\),
\begin{equation}
  K_{AB}
  =
  \chi^{-1}
  \left(A_{AB}+\frac{h_{AB}K}{4}\right),
  \qquad
  K_{ww}
  =
  \chi^{-1}
  \left(A_{ww}+\frac{h_{ww}K}{4}\right).
\end{equation}
The hidden component \(K_{ww}\) is the hidden Cartesian-like physical
extrinsic-curvature component for each suppressed direction.

The full four-dimensional conformal tracelessness relation is
\begin{equation}
  h^{AB}A_{AB}+\frac{2A_{ww}}{h_{ww}}=0.
\end{equation}
For diagonal \(h_{AB}\),
\begin{equation}
  \frac{A_{xx}}{h_{xx}}
  +
  \frac{A_{zz}}{h_{zz}}
  +
  \frac{2A_{ww}}{h_{ww}}
  =
  0,
\end{equation}
so
\begin{equation}
  A_{ww}
  =
  -\frac{h_{ww}}{2}
  \left(
    \frac{A_{xx}}{h_{xx}}
    +
    \frac{A_{zz}}{h_{zz}}
  \right).
\end{equation}

\section{Repository Checks Added After Review}

The initial Stage 3F repository script checked the \(q=h_{ww}/\chi\)
derivative identities, physical-radius identities, and the diagonal
tracelessness solve. After review, the script also checks the determinant
condition, the physical-to-conformal metric round trip, reconstruction of
\(K_{xx}\), \(K_{zz}\), and \(K_{ww}\), and the full four-dimensional
tracelessness relation with hidden multiplicity.

It also includes a guard for the dimension denominator in
\begin{equation}
  K_{ww}
  =
  \chi^{-1}
  \left(A_{ww}+\frac{h_{ww}K}{4}\right).
\end{equation}
The script reconstructs \(K_{ww}\) with a generic denominator \(p\). The
check passes for \(p=4\), but fails as expected for \(p=2\) and \(p=3\) under
a positive numeric substitution. This records that the denominator is the full
\(\mathrm{GR\_SPACEDIM}=4\), not \(\mathrm{CH\_SPACEDIM}=2\) or ordinary
\(3+1\) intuition.

\section{Limitations}

This note covers only the diagonal reduced metric. Off-diagonal \(h_{xz}\) or
\(g_{xz}\), full CCZ4 RHS terms, regularity near \(x=0\), and C++ variable
layout remain future work.

\end{document}
